\chapter{Representation Theory}

\begin{center}
    \vspace{-15pt}
    {\large\itshape "How groups act on stuff."\par}
    \vspace{-8pt}
    {\color{gray!70}\rule{0.32\textwidth}{0.5pt}\par}
\end{center}

\section{Group Actions On Sets}

\defn{Group Actions On Sets}{
    Let \(G\) be a group and \(X\) a set.

    A \textbf{left action} of \(G\) on \(X\) is a map
    \[
        \triangleright : G \times X \to X,
        \qquad
        (g,x) \mapsto g \triangleright x,
    \]
    such that
    \[
        h \triangleright (g \triangleright x) = (hg) \triangleright x,
        \qquad
        e \triangleright x = x
        \qquad
        \forall g,h \in G,\ \forall x \in X.
    \]

    A \textbf{right action} of \(G\) on \(X\) is a map
    \[
        \triangleleft : X \times G \to X,
        \qquad
        (x,g) \mapsto x \triangleleft g,
    \]
    such that
    \[
        (x \triangleleft g)\triangleleft h = x \triangleleft (gh),
        \qquad
        x \triangleleft e = x
        \qquad
        \forall g,h \in G,\ \forall x \in X.
    \]

    Equivalently, a left action is the same as a group homomorphism
    \[
        \rho : G \to \operatorname{Aut}(X),
        \qquad
        g \mapsto (x \mapsto g \triangleright x),
    \]
    and a right action is the same as an antihomomorphism
    \[
        G \to \operatorname{Aut}(X),
        \qquad
        g \mapsto (x \mapsto x \triangleleft g).
    \]
}

\corp{
    Left and right actions can be converted into each other using inverses.
    More precisely:
    \begin{itemize}
        \item if \(\triangleright\) is a left action of \(G\) on \(X\), then
        \[
            x \triangleleft g := g^{-1} \triangleright x
        \]
        defines a right action of \(G\) on \(X\);
        \item if \(\triangleleft\) is a right action of \(G\) on \(X\), then
        \[
            g \triangleright x := x \triangleleft g^{-1}
        \]
        defines a left action of \(G\) on \(X\).
    \end{itemize}
}{
    Suppose \(\triangleright\) is a left action and define
    \(x \triangleleft g := g^{-1} \triangleright x\). Then
    \[
        (x \triangleleft g)\triangleleft h
        =
        h^{-1} \triangleright (g^{-1} \triangleright x)
        =
        (h^{-1}g^{-1}) \triangleright x
        =
        (gh)^{-1} \triangleright x
        =
        x \triangleleft (gh),
    \]
    and
    \[
        x \triangleleft e = e^{-1} \triangleright x = x.
    \]
    So this is a right action. The converse is proved in the same way.
}{}

\ex{
    Examples:
    \begin{itemize}
        \item Every group \(G\) acts trivially on any set \(X\) by
        \[
            g \triangleright x = x
            \qquad \forall g \in G,\ \forall x \in X.
        \]

        \item A group \(G\) acts on itself by conjugation:
        \[
            g \triangleright h := ghg^{-1}.
        \]
        Equivalently, this is the homomorphism
        \[
            G \to \operatorname{Aut}(G),
            \qquad
            g \mapsto \operatorname{Ad}_g.
        \]

        \item If \(\Lambda \subseteq \mathbb{R}^n\) is a lattice, then its symmetry
        group \(\operatorname{Sym}(\Lambda)\) acts on \(\Lambda\).

        \item The symmetric group \(S_n\) acts on the set \(\{1,\dots,n\}\).

        \item The general linear group \(\mathrm{GL}(n,\mathbb{R})\) acts on
        \(\mathbb{R}^n\).

        \item If \(H \leq G\) is a subgroup, then \(G\) acts on the set of left
        cosets \(G/H\) by left multiplication:
        \[
            g \triangleright [g'] := [gg'].
        \]
    \end{itemize}
}

\defn{Orbit And Stabilizer}{
    Let \(\rho : G \to \operatorname{Aut}(X)\) be an action and let \(x \in X\).

    The \textbf{orbit} of \(x\), denoted by \(\operatorname{Orb}_G(x)\) or simply
    \(\mathcal{O}_x\), is the set
    \[
        \operatorname{Orb}_G(x)
        =
        \{ y \in X \mid g \triangleright x = y \text{ for some } g \in G \}
        \subseteq X.
    \]

    The \textbf{stabilizer} of \(x\), denoted by \(\operatorname{Stab}_G(x)\) or
    simply \(St_x\), is the subgroup
    \[
        \operatorname{Stab}_G(x)
        =
        \{ g \in G \mid g \triangleright x = x \}
        \leq G.
    \]
}

\ex{
        \begin{itemize}
        \item Let \(G\) act on itself by conjugation. Then
        \[
            \operatorname{Orb}_G(g)
        \]
        is the conjugacy class of \(g\), and
        \[
            \operatorname{Stab}_G(g) = \{ h \in G \mid hgh^{-1} = g \}
        \]
        is the centralizer of \(g\).

        \item The group \(\mathrm{SO}(2)\) acts on \(\mathbb{R}^2\) by rotations.
        If \(x \neq 0\), then
        \[
            \operatorname{Orb}_{\mathrm{SO}(2)}(x) \cong S^1,
            \qquad
            \operatorname{Stab}_{\mathrm{SO}(2)}(x) = \{e\}.
        \]
        For the origin one has
        \[
            \operatorname{Orb}_{\mathrm{SO}(2)}(0) = \{0\},
            \qquad
            \operatorname{Stab}_{\mathrm{SO}(2)}(0) = \mathrm{SO}(2).
        \]
    \end{itemize}
}

\thmp{Orbit-Stabilizer Theorem}{
    Let \(\rho : G \to \operatorname{Aut}(X)\) be an action of a finite group
    \(G\) on a set \(X\). Then for every \(x \in X\),
    \[
        |G| = |\operatorname{Orb}_G(x)|\,|\operatorname{Stab}_G(x)|.
    \]
}{
    Write \(St_x := \operatorname{Stab}_G(x)\) and
    \(\mathcal{O}_x := \operatorname{Orb}_G(x)\).
    Since \(St_x\) need not be normal in \(G\), the quotient \(G/St_x\) is
    considered only as the set of left cosets.

    We claim that
    \[
        \varphi : G/St_x \to \mathcal{O}_x,
        \qquad
        [g] \mapsto g \triangleright x,
    \]
    is a bijection.

    First, \(\varphi\) is well-defined. If \([g'] = [g]\), then
    \(g' = gh\) for some \(h \in St_x\). Hence
    \[
        \varphi([g'])
        =
        g' \triangleright x
        =
        (gh) \triangleright x
        =
        g \triangleright (h \triangleright x)
        =
        g \triangleright x
        =
        \varphi([g]).
    \]

    The map is surjective by the definition of the orbit. To prove injectivity,
    suppose
    \[
        g \triangleright x = g' \triangleright x.
    \]
    Applying \(g^{-1}\) to both sides gives
    \[
        (g^{-1}g') \triangleright x = x,
    \]
    so \(g^{-1}g' \in St_x\). Therefore \(g' = g(g^{-1}g')\), and hence
    \[
        [g'] = [g].
    \]
    Thus \(\varphi\) is a bijection.

    Consequently,
    \[
        |G/St_x| = |\mathcal{O}_x|.
    \]
    Since every left coset of \(St_x\) has \(|St_x|\) elements,
    \[
        |G| = |G/St_x|\,|St_x|
        =
        |\mathcal{O}_x|\,|St_x|.
    \]
}

\ex{
    \begin{enumerate}[label=\textbf{\alph*)}]
        \item Let \(G\) be a finite group and let \(H \leq G\). Then \(G\) acts
        on the set \(G/H\) of left cosets by left multiplication:
        \[
            g \triangleright [g'] := [gg'].
        \]
        The orbit of the coset \([e]\) is all of \(G/H\), and its stabilizer is
        \(H\). Indeed,
        \[
            g \triangleright [e] = [e]
            \iff
            [g] = [e]
            \iff
            g \in H.
        \]
        Hence the orbit-stabilizer theorem gives
        \[
            |G| = |G/H|\,|H|.
        \]
        In particular,
        \[
            |H| \mid |G|.
        \]

        Fix \(g \in G\) and consider the cyclic subgroup
        \[
            \langle g \rangle
            =
            \{ g^k \mid k \in \mathbb{Z} \}
            \leq G.
        \]
        Since \(G\) is finite, \(\langle g \rangle\) is finite. Its size is
        called the \textbf{order} of \(g\), denoted
        \[
            \operatorname{ord}(g) := |\langle g \rangle|.
        \]
        Equivalently, \(\operatorname{ord}(g)\) is the smallest positive
        integer \(n\) such that \(g^n = e\), and then
        \[
            \langle g \rangle = \{ e,g,g^2,\dots,g^{n-1} \},
            \qquad
            n = \operatorname{ord}(g).
        \]
        Applying the previous result to the subgroup
        \(\langle g \rangle \leq G\) gives
        \[
            \operatorname{ord}(g) \mid |G|.
        \]

        If \(|G| = p\) is prime and \(g \neq e\), then
        \[
            \operatorname{ord}(g) \neq 1
            \qquad\text{and}\qquad
            \operatorname{ord}(g) \mid p.
        \]
        Therefore \(\operatorname{ord}(g) = p\), so
        \[
            \langle g \rangle = G.
        \]
        Thus every group of prime order is cyclic, and hence
        \[
            G \cong \mathbb{Z}/p\mathbb{Z}.
        \]

        \item What are the symmetries of a cube? Let
        \[
            G
            =
            \operatorname{Sym}([-1,1]^3)
        \]
        be the full symmetry group of the cube. Let \(F\) be the set of its
        faces. Then \(G\) acts on \(F\). If \(f \in F\) is one face, then the
        orbit of \(f\) is all of \(F\), so
        \[
            \mathcal{O}_f = F,
            \qquad
            |\mathcal{O}_f| = 6.
        \]
        The stabilizer of \(f\) is the symmetry group of a square face:
        \[
            \operatorname{Stab}_G(f)
            \cong
            D_4,
            \qquad
            |D_4| = 8.
        \]
        Hence the orbit-stabilizer theorem gives
        \[
            |G|
            =
            |F|\,|D_4|
            =
            6 \cdot 8
            =
            48.
        \]

        Let \(\operatorname{Rot}\leq G\) be the subgroup of orientation
        preserving symmetries. The determinant gives a surjective homomorphism
        \[
            \det : G \to \{\pm 1\} \cong \mathbb{Z}_2
        \]
        with kernel \(\operatorname{Rot}\). Therefore
        \[
            |G| = 2|\operatorname{Rot}|,
            \qquad
            |\operatorname{Rot}| = 24.
        \]

        There are four body diagonals in the cube. If \(D\) is the set of these
        diagonals, then rotations act on \(D\), giving a homomorphism
        \[
            \operatorname{Rot}
            \to
            \operatorname{Aut}(D)
            \cong
            S_4.
        \]
        This action is faithful: a rotation fixing all four body diagonals must
        be the identity. Since both groups have order \(24\), we get
        \[
            \operatorname{Rot} \cong S_4.
        \]

        Finally, the map
        \[
            s : \mathbb{Z}_2 \to G,
            \qquad
            -1 \mapsto -I_3
        \]
        is a section of \(\det\), and \(-I_3\) commutes with every symmetry of
        the cube. Hence
        \[
            G \cong S_4 \times \mathbb{Z}_2.
        \]

        \item Let \(G\) act on itself by conjugation:
        \[
            h \triangleright g := hgh^{-1}.
        \]
        For \(g \in G\), the orbit is the conjugacy class
        \[
            \operatorname{Orb}_G(g)
            =
            \{ hgh^{-1} \mid h \in G \},
        \]
        and the stabilizer is the centralizer
        \[
            \operatorname{Stab}_G(g)
            =
            \{ h \in G \mid hgh^{-1} = g \}
            =
            \{ h \in G \mid hg = gh \}.
        \]
        Therefore the orbit-stabilizer theorem gives
        \[
            |G|
            =
            |\{ hgh^{-1} \mid h \in G \}|\,
            |\{ h \in G \mid hg = gh \}|.
        \]
    \end{enumerate}
}

\section{Representations On Vector Spaces}

\subsection{Representations And Intertwiners}

\rmkb{
    Recall that a \textbf{representation} of a group \(G\) on a vector space
    \(V\) is a group homomorphism
    \[
        \rho : G \to \operatorname{Aut}(V).
    \]
    Thus each \(g \in G\) acts on \(V\) by a structure-preserving map, namely
    an invertible linear map.
}

\defn{Maps Of Representations}{
    Let \((V,\rho)\) and \((V',\rho')\) be representations of a group \(G\).

    A \textbf{map of representations}, or \textbf{intertwiner}, is a linear map
    \[
        f : V \to V'
    \]
    such that for every \(g \in G\), the diagram
    \[
        \begin{array}{ccc}
            V & \xrightarrow{\ f\ } & V' \\
            {\scriptstyle \rho(g)}\downarrow
            &&
            \downarrow{\scriptstyle \rho'(g)} \\
            V & \xrightarrow{\ f\ } & V'
        \end{array}
    \]
    commutes. Equivalently,
    \[
        f \circ \rho(g) = \rho'(g) \circ f
        \qquad
        \forall g \in G.
    \]
}

\ex{
    \begin{itemize}
        \item Let \(V\) be a vector space. Every group \(G\) has the
        \textbf{trivial representation} on \(V\), defined by
        \[
            \rho(g) = 1_V
            \qquad
            \forall g \in G.
        \]

        \item The group \(\mathbb{Z}_2 = \{[0],[1]\}\) acts on any vector space
        \(V\) by
        \[
            \rho([0]) = 1_V,
            \qquad
            \rho([1]) = -1_V.
        \]

        \item The orthogonal group \(O(n)\) acts on \(\mathbb{R}^n\) by its
        defining representation
        \[
            \rho(A)(v) := Av.
        \]

        \item The symmetric group \(S_n\) acts on \(\mathbb{R}\) through the sign
        homomorphism:
        \[
            S_n \xrightarrow{\operatorname{sgn}} \mathbb{Z}_2
            \to
            \operatorname{Aut}(\mathbb{R}),
            \qquad
            \sigma \mapsto \bigl(x \mapsto \operatorname{sgn}(\sigma)x\bigr).
        \]

        \item The symmetric group \(S_n\) acts on \(\mathbb{R}^n\) by permuting
        the standard basis:
        \[
            \rho(\sigma)(e_i) := e_{\sigma(i)}.
        \]

        \item Let \(X\) be a set equipped with a \(G\)-action. The free vector
        space
        \[
            \mathbb{K}[X] := \operatorname{span}_{\mathbb{K}}\{x \mid x \in X\}
        \]
        carries a representation of \(G\) by linear extension:
        \[
            \rho(g)\left(\sum_i \lambda_i x_i\right)
            :=
            \sum_i \lambda_i (g \triangleright x_i).
        \]

        \item Taking \(X = G\) with the left multiplication action gives the
        \textbf{regular representation} of \(G\) on \(\mathbb{K}[G]\):
        \[
            \rho(g)(h) := gh.
        \]

        \item The symmetry group of an equilateral triangle,
        \[
            \operatorname{Sym}(\triangle) \cong S_3 \cong D_3,
        \]
        acts on \(\mathbb{R}^2\). For example, if \((12)\) is a reflection and
        \((123)\) is a rotation, then one can choose coordinates such that
        \[
            \rho((12))
            =
            \begin{pmatrix}
                -1 & 0 \\
                0 & 1
            \end{pmatrix},
            \qquad
            \rho((123))
            =
            \begin{pmatrix}
                \cos(2\pi/3) & -\sin(2\pi/3) \\
                \sin(2\pi/3) & \cos(2\pi/3)
            \end{pmatrix}.
        \]

        \item If \(\Lambda \subseteq \mathbb{R}^n\) is a lattice, then its point
        group acts on \(\mathbb{R}^n\).
    \end{itemize}
}

\subsection{Irreducibility And Schur's Lemma}

\defn{Subrepresentations And Irreducibility}{
    Let \((V,\rho)\) be a representation of a group \(G\).

    A \textbf{subrepresentation} is a linear subspace \(W \leq V\) such that
    \[
        \rho(g)(w) \in W
        \qquad
        \forall g \in G,\ \forall w \in W.
    \]
    Equivalently, \(W\) is a \(G\)-invariant subspace of \(V\).

    The representation \(V\) is \textbf{irreducible} if its only
    subrepresentations are
    \[
        \{0\}
        \qquad\text{and}\qquad
        V.
    \]
}

\propp{
    Let
    \[
        f : (V,\rho) \to (V',\rho')
    \]
    be a map of \(G\)-representations. Then
    \[
        \ker(f) \leq V
        \qquad\text{and}\qquad
        \operatorname{im}(f) \leq V'
    \]
    are subrepresentations.
}{
    Let \(v \in \ker(f)\). Then for every \(g \in G\),
    \[
        f(\rho(g)(v))
        =
        \rho'(g)(f(v))
        =
        \rho'(g)(0)
        =
        0.
    \]
    Hence \(\rho(g)(v) \in \ker(f)\), so \(\ker(f)\) is a
    subrepresentation of \(V\).

    Now let \(v' \in \operatorname{im}(f)\). Then \(v' = f(v)\) for some
    \(v \in V\). For every \(g \in G\),
    \[
        \rho'(g)(v')
        =
        \rho'(g)(f(v))
        =
        f(\rho(g)(v)).
    \]
    Thus \(\rho'(g)(v') \in \operatorname{im}(f)\), so
    \(\operatorname{im}(f)\) is a subrepresentation of \(V'\).
}

\lemp{Schur's Lemma}{
    Let \((V,\rho)\) and \((V',\rho')\) be irreducible representations of a
    group \(G\). Then every map of \(G\)-representations
    \[
        f : V \to V'
    \]
    is either the zero map or an isomorphism.

    Moreover, if \(V\) and \(V'\) are finite-dimensional complex vector spaces,
    then for any two isomorphisms of \(G\)-representations
    \[
        f,f' : V \to V'
    \]
    there exists \(\lambda \in \mathbb{C}\) such that
    \[
        f' = \lambda f.
    \]
}{
    Let \(f : V \to V'\) be a map of \(G\)-representations. By the previous
    proposition, \(\ker(f)\) is a subrepresentation of \(V\), and
    \(\operatorname{im}(f)\) is a subrepresentation of \(V'\). Since \(V\) and
    \(V'\) are irreducible, we have
    \[
        \ker(f) \in \{\{0\},V\},
        \qquad
        \operatorname{im}(f) \in \{\{0\},V'\}.
    \]
    If \(\operatorname{im}(f)=\{0\}\), then \(f=0\). Otherwise
    \(\operatorname{im}(f)=V'\). In this case \(\ker(f) \neq V\), so
    \(\ker(f)=\{0\}\). Hence \(f\) is both surjective and injective, and
    therefore an isomorphism.

    Now assume that \(V\) and \(V'\) are finite-dimensional complex vector
    spaces, and let
    \[
        f,f' : V \to V'
    \]
    be isomorphisms of \(G\)-representations. Then
    \[
        f^{-1} \circ f' : V \to V
    \]
    is a map of \(G\)-representations. Since \(V\) is a finite-dimensional
    complex vector space, \(f^{-1}\circ f'\) has an eigenvalue
    \(\lambda \in \mathbb{C}\). Thus
    \[
        f^{-1}\circ f' - \lambda 1_V
    \]
    is not injective. By the first part, it cannot be an isomorphism, so it
    must be the zero map. Therefore
    \[
        f^{-1}\circ f' = \lambda 1_V,
    \]
    and multiplying by \(f\) gives
    \[
        f' = \lambda f.
    \]
}

\ex{
    The second part of Schur's lemma need not hold over \(\mathbb{R}\).
    Consider the representation of \(\mathbb{Z}_4\) on \(\mathbb{R}^2\) by
    rotations:
    \[
        \rho([n])
        =
        \begin{pmatrix}
            0 & 1 \\
            -1 & 0
        \end{pmatrix}^{n}.
    \]
    Then both
    \[
        1_{\mathbb{R}^2}
        \qquad\text{and}\qquad
        \begin{pmatrix}
            0 & 1 \\
            -1 & 0
        \end{pmatrix}
        :
        \mathbb{R}^2 \to \mathbb{R}^2
    \]
    are maps of \(\mathbb{Z}_4\)-representations. Indeed, both commute with
    the action of the generator \([1]\), and hence with the action of every
    element of \(\mathbb{Z}_4\).

    However, these two maps are not scalar multiples of each other over
    \(\mathbb{R}\), since the rotation matrix is not of the form
    \(\lambda 1_{\mathbb{R}^2}\) for any
    \(\lambda \in \mathbb{R}\). Thus the complex-vector-space assumption in
    Schur's lemma is essential for the scalar uniqueness statement.
}

\rmkb{
    From now on, unless stated otherwise, representations will be taken over
    complex vector spaces.
}

\subsection{Central Elements And Abelian Groups}

\corp{
    Let \((V,\rho)\) be a finite-dimensional irreducible representation of a
    group \(G\). Then every element of the center
    \[
        Z(G) := \{ c \in G \mid cg = gc \text{ for all } g \in G \}
    \]
    acts by a scalar. In other words, for every \(c \in Z(G)\) there exists
    \(\lambda_c \in \mathbb{C}\) such that
    \[
        \rho(c) = \lambda_c 1_V.
    \]
}{
    Let \(c \in Z(G)\). Then for every \(g \in G\),
    \[
        \rho(c) \circ \rho(g)
        =
        \rho(cg)
        =
        \rho(gc)
        =
        \rho(g) \circ \rho(c).
    \]
    Thus \(\rho(c) : V \to V\) is a map of \(G\)-representations. Since \(V\)
    is a finite-dimensional complex vector space, \(\rho(c)\) has an
    eigenvalue \(\lambda_c \in \mathbb{C}\). Then
    \[
        \rho(c) - \lambda_c 1_V
    \]
    is again a map of \(G\)-representations, but it is not injective. By
    Schur's lemma it cannot be an isomorphism, so it must be the zero map.
    Hence
    \[
        \rho(c) = \lambda_c 1_V.
    \]
}{}

\ex{
    Let \(G\) be an abelian group. Then
    \[
        Z(G) = G.
    \]
    Hence, if \((V,\rho)\) is a finite-dimensional irreducible representation
    of \(G\), every element of \(G\) acts by a scalar.

    Therefore every linear subspace of \(V\) is preserved by the action of
    \(G\). Since \(V\) is irreducible, this is only possible when
    \[
        \dim V = 1.
    \]
    Thus such representations are one-dimensional. Equivalently, they are
    homomorphisms
    \[
        \rho : G \to \mathbb{C}^{\times}.
    \]

    For example, let \(G = \mathbb{Z}_n\). If \([1]\) is the generator, then
    \[
        \rho([1])^n = \rho([0]) = 1.
    \]
    Hence \(\rho([1])\) must be an \(n\)-th root of unity. For
    \(k \in \{0,1,\dots,n-1\}\), define
    \[
        \rho_k([1])
        =
        e^{2\pi i k/n}.
    \]
    Then
    \[
        \rho_k([a])
        =
        e^{2\pi i k a/n}
    \]
    defines a one-dimensional irreducible representation of \(\mathbb{Z}_n\),
    and these are all of them.
}

\subsection{Decomposition Of Representations}

\defn{Constructions From Old Representations}{
    Let \(G\) be a group, and let \((V_1,\rho_1)\) and \((V_2,\rho_2)\) be
    representations of \(G\).

    \begin{itemize}
        \item The \textbf{direct sum representation} is
        \[
            (V_1 \oplus V_2,\rho_1 \oplus \rho_2),
        \]
        where
        \[
            (\rho_1 \oplus \rho_2)(g)(v_1,v_2)
            :=
            (\rho_1(g)v_1,\rho_2(g)v_2).
        \]

        \item The \textbf{tensor product representation} is
        \[
            (V_1 \otimes V_2,\rho_1 \otimes \rho_2),
        \]
        where on pure tensors
        \[
            (\rho_1 \otimes \rho_2)(g)(v_1 \otimes v_2)
            :=
            \rho_1(g)v_1 \otimes \rho_2(g)v_2,
        \]
        and then extended linearly.

        \item The \textbf{dual representation} of \((V_1,\rho_1)\) is
        \[
            (V_1^*,\rho_1^*),
            \qquad
            V_1^* := \operatorname{Hom}(V_1,\mathbb{C}),
        \]
        where
        \[
            (\rho_1^*(g)f)(v)
            :=
            f\bigl(\rho_1(g^{-1})v\bigr).
        \]

        \item The \textbf{Hom representation} is the representation on
        \(\operatorname{Hom}(V_1,V_2)\) defined by
        \[
            \rho(g) : \operatorname{Hom}(V_1,V_2)
            \to
            \operatorname{Hom}(V_1,V_2),
        \]
        with
        \[
            (\rho(g)F)(v)
            :=
            \rho_2(g)\Bigl(F\bigl(\rho_1(g^{-1})v\bigr)\Bigr).
        \]
    \end{itemize}
}

\propp{
    The direct sum, tensor product, dual, and Hom constructions above define
    representations of \(G\).
}{
    For the direct sum, one has
    \[
        (\rho_1 \oplus \rho_2)(e)(v_1,v_2)
        =
        (v_1,v_2),
    \]
    and
    \[
        (\rho_1 \oplus \rho_2)(gh)(v_1,v_2)
        =
        (\rho_1(g)\rho_1(h)v_1,\rho_2(g)\rho_2(h)v_2)
        =
        ((\rho_1 \oplus \rho_2)(g)(\rho_1 \oplus \rho_2)(h))(v_1,v_2).
    \]

    For the tensor product, the identity element acts as the identity. Moreover,
    on pure tensors,
    \[
        (\rho_1 \otimes \rho_2)(gh)(v_1 \otimes v_2)
        =
        \rho_1(g)\rho_1(h)v_1 \otimes \rho_2(g)\rho_2(h)v_2,
    \]
    which is equal to
    \[
        ((\rho_1 \otimes \rho_2)(g)(\rho_1 \otimes \rho_2)(h))
        (v_1 \otimes v_2).
    \]
    Since pure tensors span \(V_1 \otimes V_2\), this proves the representation
    identity on the whole tensor product.

    For the dual representation, let \(f \in V_1^*\) and \(v \in V_1\). Then
    \[
        (\rho_1^*(e)f)(v)
        =
        f(\rho_1(e^{-1})v)
        =
        f(v),
    \]
    and
    \[
        (\rho_1^*(gh)f)(v)
        =
        f(\rho_1((gh)^{-1})v)
        =
        f(\rho_1(h^{-1})\rho_1(g^{-1})v).
    \]
    On the other hand,
    \[
        (\rho_1^*(g)(\rho_1^*(h)f))(v)
        =
        (\rho_1^*(h)f)(\rho_1(g^{-1})v)
        =
        f(\rho_1(h^{-1})\rho_1(g^{-1})v).
    \]
    Hence \(\rho_1^*(gh)=\rho_1^*(g)\rho_1^*(h)\).

    Finally, for the Hom representation, the identity element acts as the
    identity. For \(F \in \operatorname{Hom}(V_1,V_2)\) and \(v \in V_1\), we
    compute
    \[
        (\rho(gh)F)(v)
        =
        \rho_2(gh)F(\rho_1((gh)^{-1})v)
        =
        \rho_2(g)\rho_2(h)F(\rho_1(h^{-1})\rho_1(g^{-1})v).
    \]
    Also,
    \[
        (\rho(g)(\rho(h)F))(v)
        =
        \rho_2(g)\Bigl((\rho(h)F)(\rho_1(g^{-1})v)\Bigr)
        =
        \rho_2(g)\rho_2(h)F(\rho_1(h^{-1})\rho_1(g^{-1})v).
    \]
    Thus \(\rho(gh)=\rho(g)\rho(h)\), so this is a representation.
}

\rmkb{
    Let \(V_1\) and \(V_2\) be finite-dimensional representations of \(G\).
    Then there is a natural linear isomorphism
    \[
        \Phi : V_2 \otimes V_1^* \to \operatorname{Hom}(V_1,V_2),
        \qquad
        v_2 \otimes f \mapsto \bigl(v_1 \mapsto f(v_1)v_2\bigr).
    \]
    This is an isomorphism of representations. Indeed, for \(g \in G\),
    \(v_1 \in V_1\), \(v_2 \in V_2\), and \(f \in V_1^*\), we have
    \[
        \Phi\bigl((\rho_2 \otimes \rho_1^*)(g)(v_2 \otimes f)\bigr)(v_1)
        =
        f(\rho_1(g^{-1})v_1)\rho_2(g)v_2.
    \]
    On the other hand,
    \[
        (\rho(g)\Phi(v_2 \otimes f))(v_1)
        =
        \rho_2(g)\Bigl(\Phi(v_2 \otimes f)(\rho_1(g^{-1})v_1)\Bigr)
        =
        f(\rho_1(g^{-1})v_1)\rho_2(g)v_2.
    \]
    Hence \(\Phi\) intertwines the two representations.
}

\defn{Indecomposable Representation}{
    A representation \((V,\rho)\) of a group \(G\) is
    \textbf{indecomposable} if it cannot be written as a direct sum of two
    nonzero subrepresentations. In other words, there do not exist nonzero
    subrepresentations \(W_1,W_2 \leq V\) such that
    \[
        V = W_1 \oplus W_2.
    \]
}

\rmkb{
    Every irreducible representation is indecomposable: if
    \(V = W_1 \oplus W_2\) with \(W_1,W_2\) nonzero subrepresentations, then
    \(W_1\) and \(W_2\) are proper nonzero subrepresentations of \(V\), so
    \(V\) is not irreducible.

    The converse is not true in general: an indecomposable representation can
    still have nonzero proper subrepresentations.
}

\ex{
    \begin{itemize}
        \item Every one-dimensional representation is indecomposable, since a
        one-dimensional vector space cannot be written as a direct sum of two
        nonzero subspaces.

        \item Consider the representation of \(\mathbb{Z}\) on
        \(\mathbb{C}^2\) defined by
        \[
            \rho(1)
            =
            \begin{pmatrix}
                1 & 1 \\
                0 & 1
            \end{pmatrix}.
        \]
        Then the line
        \[
            W = \operatorname{span}\left\{
                \begin{pmatrix}
                    1 \\
                    0
                \end{pmatrix}
            \right\}
            \subseteq \mathbb{C}^2
        \]
        is a subrepresentation, so the representation is reducible.

        However, it is indecomposable. Indeed, a decomposition into two
        one-dimensional subrepresentations would give a basis of
        \(\mathbb{C}^2\) consisting of eigenvectors of \(\rho(1)\). This would
        make \(\rho(1)\) diagonalizable. But
        \[
            \begin{pmatrix}
                1 & 1 \\
                0 & 1
            \end{pmatrix}
        \]
        is not diagonalizable; its only eigenspace is \(W\). Therefore this
        representation is reducible but indecomposable.
    \end{itemize}
}

\defn{Semisimple Representations}{
    Let \((V,\rho)\) be a representation of a group \(G\). We say that \(V\) is
    \textbf{semisimple}, or \textbf{completely reducible}, if every
    subrepresentation \(W \leq V\) has a complementary subrepresentation.
    That is, for every subrepresentation \(W \leq V\), there exists a
    subrepresentation \(W' \leq V\) such that
    \[
        V = W \oplus W'.
    \]

    If \(V\) is finite-dimensional, this is equivalent to saying that
    \[
        V \cong \bigoplus_i W_i
    \]
    with each \(W_i\) irreducible.
}

\ex{
    Recall that
    \[
        \mathrm{SU}(2)
        =
        \{ U \in \operatorname{Mat}(2 \times 2,\mathbb{C})
        \mid U^\dagger U = 1,\ \det(U)=1 \}
    \]
    has a defining two-dimensional representation \(V_{1/2} = \mathbb{C}^2\).
    For this representation one has
    \[
        V_{1/2} \otimes V_{1/2}
        \cong
        \operatorname{Sym}^2(V_{1/2})
        \oplus
        \bigwedge^2 V_{1/2}
    \]
    as representations.

    More generally, the irreducible representations of \(\mathrm{SU}(2)\) are
    indexed by
    \[
        j \in \frac{1}{2}\mathbb{N}_0,
    \]
    and are given by
    \[
        V_j := \operatorname{Sym}^{2j}(V_{1/2}).
    \]
    The tensor product decomposes according to the Clebsch-Gordan rule
    \[
        V_k \otimes V_\ell
        \cong
        \bigoplus_{j=|k-\ell|}^{k+\ell} V_j,
    \]
    where \(j\) increases in steps of \(1\). This is the representation
    theoretic form of addition of angular momentum in quantum mechanics.

    We will come back to this example later when discussing Lie groups.
}

\rmkb{
    We want to show that, for finite groups, irreducible and indecomposable
    representations are the same thing. Equivalently, every finite-dimensional
    representation of a finite group should split as a direct sum of
    irreducible representations.

    The subtle point is that a subspace does not have a preferred complement
    in general. For Hilbert spaces, however, there is a natural candidate: the
    orthogonal complement. This motivates the following definition.
}

\defn{Unitary Representation}{
    A \textbf{unitary representation} of a group \(G\) is a Hilbert space
    \(\mathcal{H}\) together with a group homomorphism
    \[
        \rho : G \to U(\mathcal{H}),
    \]
    where \(U(\mathcal{H})\) denotes the group of unitary automorphisms of
    \(\mathcal{H}\).
}

\propp{
    Let \((V,\rho)\) be a finite-dimensional representation of a finite group
    \(G\). Then there exists a positive Hermitian inner product on \(V\) such
    that \(\rho\) is unitary. In other words, every finite-dimensional
    representation of a finite group can be lifted to a unitary
    representation.
}{
    Let
    \[
        \langle -,- \rangle
    \]
    be any positive Hermitian inner product on \(V\). Define a new inner
    product by averaging over the group:
    \[
        \langle v,w\rangle_G
        :=
        \sum_{g \in G}
        \langle \rho(g)v,\rho(g)w\rangle.
    \]
    Since \(G\) is finite, this sum is finite.

    The form \(\langle -,-\rangle_G\) is again Hermitian and positive. Indeed,
    each summand is Hermitian, and for \(v \neq 0\) one has
    \[
        \langle v,v\rangle_G
        =
        \sum_{g \in G}
        \langle \rho(g)v,\rho(g)v\rangle
        >
        0,
    \]
    because every \(\rho(g)\) is invertible, so \(\rho(g)v \neq 0\).

    It remains to check that the representation is unitary with respect to
    this new inner product. Let \(h \in G\). Then
    \[
        \langle \rho(h)v,\rho(h)w\rangle_G
        =
        \sum_{g \in G}
        \langle \rho(g)\rho(h)v,\rho(g)\rho(h)w\rangle
    \]
    and hence
    \[
        \langle \rho(h)v,\rho(h)w\rangle_G
        =
        \sum_{g \in G}
        \langle \rho(gh)v,\rho(gh)w\rangle.
    \]
    As \(g\) runs through \(G\), so does \(gh\). Therefore, after the change of
    variables \(g' = gh\),
    \[
        \langle \rho(h)v,\rho(h)w\rangle_G
        =
        \sum_{g' \in G}
        \langle \rho(g')v,\rho(g')w\rangle
        =
        \langle v,w\rangle_G.
    \]
    Thus every \(\rho(h)\) preserves \(\langle -,-\rangle_G\), so \(\rho\) is
    unitary.
}

\thmp{Maschke's Theorem}{
    Every finite-dimensional representation of a finite group is completely
    reducible.
}{
    Let \((V,\rho)\) be a finite-dimensional representation of a finite group
    \(G\). By the previous proposition, we may choose a positive Hermitian
    inner product
    \[
        \langle -,-\rangle_G
    \]
    on \(V\) such that \(\rho\) is unitary.

    Let \(W \leq V\) be a subrepresentation. Define its orthogonal complement
    with respect to this inner product by
    \[
        W^\perp
        :=
        \{ v \in V \mid \langle v,w\rangle_G = 0
        \text{ for all } w \in W \}.
    \]
    Since \(V\) is finite-dimensional, we have a direct sum decomposition of
    vector spaces
    \[
        V = W \oplus W^\perp.
    \]

    It remains to show that \(W^\perp\) is again a subrepresentation. Let
    \(v \in W^\perp\), \(g \in G\), and \(w \in W\). Since \(\rho(g)\) is
    unitary,
    \[
        \langle \rho(g)v,w\rangle_G
        =
        \langle \rho(g)v,\rho(g)\rho(g^{-1})w\rangle_G
        =
        \langle v,\rho(g^{-1})w\rangle_G.
    \]
    Because \(W\) is a subrepresentation, \(\rho(g^{-1})w \in W\). Since
    \(v \in W^\perp\), this gives
    \[
        \langle \rho(g)v,w\rangle_G = 0.
    \]
    Hence \(\rho(g)v \in W^\perp\), so \(W^\perp\) is a subrepresentation.

    Thus every subrepresentation \(W \leq V\) has a complementary
    subrepresentation \(W^\perp\). Therefore \(V\) is completely reducible.
}

\corp{
    Let \(G\) be a finite group. A finite-dimensional representation \(V\) of
    \(G\) is indecomposable if and only if it is irreducible.
    Equivalently, every finite-dimensional representation of \(G\) can be
    written as
    \[
        V \cong \bigoplus_i N_i V_i,
    \]
    where each \(V_i\) is irreducible and \(N_i \geq 0\) is the multiplicity of
    \(V_i\) in \(V\). Here
    \[
        N_i V_i := \underbrace{V_i \oplus \cdots \oplus V_i}_{N_i\text{ times}}.
    \]
}{
    By Maschke's theorem, every finite-dimensional representation of \(G\) is
    completely reducible.

    We first prove the direct sum decomposition. If \(V=0\), there is nothing
    to show. Otherwise, choose a nonzero subrepresentation \(W \leq V\) of
    minimal dimension. Then \(W\) is irreducible: any nonzero proper
    subrepresentation of \(W\) would contradict the minimality of \(W\).
    By complete reducibility, there is a subrepresentation \(W' \leq V\) such
    that
    \[
        V = W \oplus W'.
    \]
    Repeating this argument for \(W'\), and using finite-dimensionality, the
    process terminates and gives a decomposition of \(V\) into irreducible
    subrepresentations. Grouping isomorphic irreducible summands gives
    \[
        V \cong \bigoplus_i N_i V_i.
    \]

    Now suppose that \(V\) is indecomposable. In such a decomposition, there
    can be only one irreducible summand and it can occur only once; otherwise
    \(V\) would split as a direct sum of two nonzero subrepresentations. Hence
    \(V\) is irreducible.

    Conversely, every irreducible representation is indecomposable, as noted
    above. Therefore indecomposable and irreducible representations agree for
    finite groups.
}{}

\ex{
    Let
    \[
        \mathbb{Z}_n = \mathbb{Z}/n\mathbb{Z}
    \]
    be the cyclic group, thought of as \(n\) equally spaced points on a circle.
    Consider the vector space of complex-valued functions
    \[
        \mathcal{O}(\mathbb{Z}_n)
        :=
        \{ f : \mathbb{Z}_n \to \mathbb{C} \}.
    \]
    The group \(\mathbb{Z}_n\) acts on this space by translations:
    \[
        (\rho([a])f)([x])
        :=
        f([x+a]).
    \]

    For \(k=0,1,\dots,n-1\), define
    \[
        f_k([x])
        :=
        e^{2\pi i kx/n}.
    \]
    These functions are eigenvectors for every translation, since
    \[
        (\rho([a])f_k)([x])
        =
        f_k([x+a])
        =
        e^{2\pi i k(x+a)/n}
        =
        e^{2\pi i ka/n} f_k([x]).
    \]
    Hence the line
    \[
        \mathbb{C}_k := \mathbb{C} f_k
    \]
    is a one-dimensional subrepresentation of \(\mathcal{O}(\mathbb{Z}_n)\).
    Since the functions \(f_0,\dots,f_{n-1}\) form a basis of
    \(\mathcal{O}(\mathbb{Z}_n)\), we get the decomposition
    \[
        \mathcal{O}(\mathbb{Z}_n)
        =
        \bigoplus_{k=0}^{n-1} \mathbb{C}_k
        =
        \bigoplus_{k=0}^{n-1} \mathbb{C} f_k.
    \]
    Thus the decomposition of the regular representation of
    \(\mathbb{Z}_n\) is exactly the discrete Fourier transform.
}

\subsection{Character Theory}

\defn{Characters And Class Functions}{
    Let \((V,\rho)\) be a finite-dimensional representation of a group \(G\).
    The \textbf{character} of \(V\) is the function
    \[
        \chi_V : G \to \mathbb{C},
        \qquad
        g \mapsto \operatorname{tr}_V(\rho(g)).
    \]

    A function \(f : G \to \mathbb{C}\) is called a \textbf{class function} if
    it is invariant under conjugation, that is,
    \[
        f(hgh^{-1}) = f(g)
        \qquad
        \forall g,h \in G.
    \]
}

\propp{
    Characters are class functions.
}{
    Let \((V,\rho)\) be a finite-dimensional representation of \(G\). Then for
    every \(g,h \in G\),
    \[
        \chi_V(hgh^{-1})
        =
        \operatorname{tr}_V(\rho(hgh^{-1}))
        =
        \operatorname{tr}_V(\rho(h)\rho(g)\rho(h^{-1})).
    \]
    Using cyclicity of the trace, this becomes
    \[
        \operatorname{tr}_V(\rho(h)\rho(g)\rho(h^{-1}))
        =
        \operatorname{tr}_V(\rho(g)\rho(h^{-1})\rho(h))
        =
        \operatorname{tr}_V(\rho(g)).
    \]
    Hence
    \[
        \chi_V(hgh^{-1}) = \chi_V(g),
    \]
    so \(\chi_V\) is invariant under conjugation.
}

\rmkb{
    Since characters are class functions, it is enough to record their values
    on the conjugacy classes of \(G\). A table whose rows are irreducible
    characters and whose columns are conjugacy classes is called a
    \textbf{character table}.
}

\ex{
    The group \(S_3\) has three conjugacy classes, represented by
    \[
        1,
        \qquad
        (12),
        \qquad
        (123).
    \]
    It has three irreducible representations: the trivial representation, the
    alternating representation, and the standard representation. Their
    character table is
    \[
        \begin{array}{c|ccc}
            S_3 & 1 & (12) & (123) \\
            \hline
            \chi_{\mathrm{triv}} & 1 & 1 & 1 \\
            \chi_{\mathrm{alt}}  & 1 & -1 & 1 \\
            \chi_{\mathrm{std}}  & 2 & 0 & -1
        \end{array}
    \]
}

\ex{
    For \(\mathbb{Z}_3\), every conjugacy class consists of a single element.
    The irreducible characters are indexed by \(k=0,1,2\). Their character
    table is
    \[
        \begin{array}{c|ccc}
            \mathbb{Z}_3 & [0] & [1] & [2] \\
            \hline
            k=0 & 1 & 1 & 1 \\
            k=1 & 1 & e^{2\pi i/3} & e^{4\pi i/3} \\
            k=2 & 1 & e^{4\pi i/3} & e^{8\pi i/3}
        \end{array}
    \]
}

\rmkb{
    Characters behave well with respect to direct sums and tensor products:
    \[
        \chi_{V \oplus W} = \chi_V + \chi_W,
        \qquad
        \chi_{V \otimes W} = \chi_V \chi_W.
    \]
    The second identity follows from the fact that the trace of a tensor
    product of linear maps is the product of the traces.
}

\ex{
    For the standard representation of \(S_3\), the character is
    \[
        \chi_{\mathrm{std}} = (2,0,-1)
    \]
    on the conjugacy classes represented by \(1,(12),(123)\). Therefore
    \[
        \chi_{\mathrm{std} \otimes \mathrm{std}}
        =
        \chi_{\mathrm{std}}\chi_{\mathrm{std}}
        =
        (4,0,1).
    \]
    On the other hand,
    \[
        \chi_{\mathrm{std}}
        +
        \chi_{\mathrm{alt}}
        +
        \chi_{\mathrm{triv}}
        =
        (2,0,-1) + (1,-1,1) + (1,1,1)
        =
        (4,0,1).
    \]
    Hence the character table suggests the decomposition
    \[
        \mathrm{Std} \otimes \mathrm{Std}
        \cong
        \mathrm{Std} \oplus \mathrm{Alt} \oplus \mathrm{Triv}.
    \]
    In particular, the characters in the table are linearly independent.
}

\defn{Invariant Subspace}{
    Let \(G\) be a group and let \((V,\rho)\) be a representation of \(G\).
    The \textbf{subrepresentation of \(G\)-invariants} is
    \[
        V^G
        :=
        \{ v \in V \mid \rho(g)(v) = v \text{ for all } g \in G \}
        \leq V.
    \]
}

\rmkb{
    The subspace \(V^G\) is the part of \(V\) on which \(G\) acts trivially.
    Thus, for finite-dimensional representations of finite groups, one can
    write
    \[
        V \cong \dim(V^G)\,\mathrm{Triv} \oplus V',
    \]
    where \(V'\) contains no copy of the trivial representation.

    In particular, if \(V_i\) is irreducible but not trivial, then
    \[
        V_i^G = \{0\}.
    \]
}

\propp{
    Let \(G\) be a finite group and let \((V,\rho)\) be a representation of
    \(G\). The map
    \[
        N_G : V \to V,
        \qquad
        v \mapsto \frac{1}{|G|}\sum_{g \in G}\rho(g)(v)
    \]
    is a projection operator onto \(V^G\). In other words,
    \[
        N_G \circ N_G = N_G,
        \qquad
        \operatorname{im}(N_G) = V^G.
    \]
}{
    First we check that \(N_G\) intertwines the \(G\)-action. Let \(g \in G\)
    and \(v \in V\). Then
    \[
        N_G(\rho(g)(v))
        =
        \frac{1}{|G|}\sum_{h \in G}\rho(hg)(v).
    \]
    As \(h\) runs through \(G\), so does \(h' = hg\). Therefore
    \[
        N_G(\rho(g)(v))
        =
        \frac{1}{|G|}\sum_{h' \in G}\rho(h')(v)
        =
        N_G(v).
    \]
    Also,
    \[
        \rho(g)(N_G(v))
        =
        \rho(g)\left(\frac{1}{|G|}\sum_{h \in G}\rho(h)(v)\right)
        =
        \frac{1}{|G|}\sum_{h \in G}\rho(gh)(v).
    \]
    Again, as \(h\) runs through \(G\), so does \(h' = gh\), and hence
    \[
        \rho(g)(N_G(v)) = N_G(v).
    \]
    Thus \(\operatorname{im}(N_G) \subseteq V^G\).

    Conversely, if \(v \in V^G\), then \(\rho(g)(v)=v\) for every \(g \in G\).
    Hence
    \[
        N_G(v)
        =
        \frac{1}{|G|}\sum_{g \in G}\rho(g)(v)
        =
        \frac{1}{|G|}\sum_{g \in G}v
        =
        v.
    \]
    Therefore \(V^G \subseteq \operatorname{im}(N_G)\), and so
    \[
        \operatorname{im}(N_G)=V^G.
    \]

    Finally, since \(N_G(v)\in V^G\) for every \(v\in V\), the previous
    computation gives
    \[
        N_G(N_G(v)) = N_G(v).
    \]
    Thus \(N_G^2=N_G\), so \(N_G\) is a projection operator onto \(V^G\).
}

\corp{
    Let \(G\) be a finite group and let \(V\) be a finite-dimensional
    representation of \(G\). Then
    \[
        \dim(V^G)
        =
        \frac{1}{|G|}\sum_{g \in G}\chi_V(g).
    \]
    In particular, if \(V\) is irreducible and not trivial, then
    \[
        \sum_{g \in G}\chi_V(g) = 0.
    \]
}{
    By the previous proposition, \(N_G\) is a projection onto \(V^G\). Thus,
    with respect to a decomposition
    \[
        V \cong V^G \oplus \ker(N_G),
    \]
    the operator \(N_G\) has matrix
    \[
        \begin{pmatrix}
            I_{\dim(V^G)} & 0 \\
            0 & 0
        \end{pmatrix}.
    \]
    Therefore
    \[
        \operatorname{tr}(N_G)=\dim(V^G).
    \]
    On the other hand,
    \[
        \operatorname{tr}(N_G)
        =
        \operatorname{tr}\left(
            \frac{1}{|G|}\sum_{g \in G}\rho(g)
        \right)
        =
        \frac{1}{|G|}\sum_{g \in G}\operatorname{tr}(\rho(g))
        =
        \frac{1}{|G|}\sum_{g \in G}\chi_V(g).
    \]
    This proves the formula. If \(V\) is irreducible and not trivial, then
    \(V^G=\{0\}\), and hence the sum of the character values is zero.
}{}

\rmkb{
    If \(V\) and \(W\) are representations of \(G\), then
    \(\operatorname{Hom}(V,W)\) is again a representation. Its invariant
    subspace is
    \[
        \operatorname{Hom}(V,W)^G
        =
        \{ f : V \to W \mid
        f(\rho_V(g)v)=\rho_W(g)f(v)
        \text{ for all } g \in G,\ v \in V \}.
    \]
    Thus \(\operatorname{Hom}(V,W)^G\) is precisely the space of maps of
    \(G\)-representations from \(V\) to \(W\).
}

\corp{
    Let \(V\) and \(W\) be irreducible finite-dimensional complex
    representations of a finite group \(G\). Then
    \[
        \frac{1}{|G|}\sum_{g \in G}
        \overline{\chi_V(g)}\,\chi_W(g)
        =
        \begin{cases}
            1 & \text{if } V \cong W, \\
            0 & \text{otherwise}.
        \end{cases}
    \]
}{
    By Schur's lemma,
    \[
        \operatorname{Hom}(V,W)^G
        =
        \begin{cases}
            \mathbb{C} & \text{if } V \cong W, \\
            \{0\} & \text{otherwise}.
        \end{cases}
    \]
    Hence
    \[
        \dim \operatorname{Hom}(V,W)^G
        =
        \begin{cases}
            1 & \text{if } V \cong W, \\
            0 & \text{otherwise}.
        \end{cases}
    \]

    Applying the previous corollary to the representation
    \(\operatorname{Hom}(V,W)\), we get
    \[
        \dim \operatorname{Hom}(V,W)^G
        =
        \frac{1}{|G|}\sum_{g \in G}
        \chi_{\operatorname{Hom}(V,W)}(g).
    \]
    Since
    \[
        \operatorname{Hom}(V,W) \cong V^* \otimes W,
    \]
    its character is
    \[
        \chi_{\operatorname{Hom}(V,W)}(g)
        =
        \chi_{V^*}(g)\chi_W(g)
        =
        \overline{\chi_V(g)}\,\chi_W(g).
    \]
    Therefore
    \[
        \frac{1}{|G|}\sum_{g \in G}
        \overline{\chi_V(g)}\,\chi_W(g)
        =
        \begin{cases}
            1 & \text{if } V \cong W, \\
            0 & \text{otherwise}.
        \end{cases}
    \]
}{}

\corp{
    The characters of irreducible finite-dimensional complex representations
    of a finite group \(G\) are orthonormal with respect to the scalar product
    \[
        (F,H)
        :=
        \frac{1}{|G|}\sum_{g \in G}\overline{F(g)}H(g)
    \]
    on the space of class functions on \(G\).
}{
    This is exactly the previous corollary, written in the notation
    \[
        (\chi_V,\chi_W)
        =
        \frac{1}{|G|}\sum_{g \in G}
        \overline{\chi_V(g)}\,\chi_W(g).
    \]
}{}

\corp{
    Every finite-dimensional representation \(V\) of a finite group \(G\) is
    completely determined, up to isomorphism, by its character \(\chi_V\).
    More precisely, if
    \[
        V \cong \bigoplus_i N_i V_i
    \]
    is a decomposition into irreducible representations, then
    \[
        N_i = (\chi_{V_i},\chi_V).
    \]
}{
    By Maschke's theorem, write
    \[
        V \cong \bigoplus_i N_i V_i
    \]
    with the \(V_i\) irreducible. Then
    \[
        \chi_V = \sum_i N_i \chi_{V_i}.
    \]
    Taking the scalar product with \(\chi_{V_j}\), and using the
    orthonormality of irreducible characters, gives
    \[
        (\chi_{V_j},\chi_V)
        =
        \sum_i N_i(\chi_{V_j},\chi_{V_i})
        =
        N_j.
    \]
    Thus the character determines all multiplicities \(N_i\), and hence
    determines \(V\) up to isomorphism.
}{}

\rmkb{
    Since irreducible characters are linearly independent, there can be only
    finitely many irreducible representations of a finite group up to
    isomorphism.
}

\ex{
    Let
    \[
        D_4 = \langle r,s \mid r^4=s^2=1,\ srs=r^{-1}\rangle
    \]
    be the symmetry group of the square, where \(r\) is rotation by
    \(90^\circ\), and \(s\) is a reflection. Its conjugacy classes are
    \[
        \{1\},
        \qquad
        \{r,r^3\},
        \qquad
        \{r^2\},
        \qquad
        \{s,sr^2\},
        \qquad
        \{sr,sr^3\}.
    \]

    The one-dimensional representations are homomorphisms
    \[
        \varphi : D_4 \to \mathbb{C}^{\times}.
    \]
    Since
    \[
        \varphi(r)
        =
        \varphi(srs)
        =
        \varphi(s)\varphi(r)\varphi(s)
        =
        \varphi(r)^{-1},
    \]
    we must have \(\varphi(r)=\pm 1\). Also, since \(s^2=1\), one has
    \(\varphi(s)=\pm 1\). Thus the one-dimensional representations are
    labelled by the choices
    \[
        (\varphi(r),\varphi(s))
        \in
        \{(1,1),(1,-1),(-1,1),(-1,-1)\}.
    \]

    There is also the standard two-dimensional representation
    \[
        D_4 \to \operatorname{GL}(\mathbb{R}^2),
    \]
    coming from the action of \(D_4\) on the square in the plane. Its character
    is
    \[
        \chi_{\mathrm{std}} = (2,0,-2,0,0)
    \]
    on the conjugacy classes listed above. Since
    \[
        (\chi_{\mathrm{std}},\chi_{\mathrm{std}})
        =
        \frac{1}{8}(2^2+(-2)^2)
        =
        1,
    \]
    the standard representation is irreducible.

    The character table is
    \[
        \begin{array}{c|ccccc}
            D_4 & 1 & r & r^2 & s & sr \\
            \hline
            (1,1)     & 1 & 1  & 1  & 1  & 1  \\
            (-1,1)    & 1 & -1 & 1  & 1  & -1 \\
            (1,-1)    & 1 & 1  & 1  & -1 & -1 \\
            (-1,-1)   & 1 & -1 & 1  & -1 & 1  \\
            \mathrm{std} & 2 & 0  & -2 & 0  & 0
        \end{array}
    \]
    Here the column labels stand for the conjugacy classes represented by
    those elements.

    Now consider \(\mathrm{std}\otimes \mathrm{std}\). Its character is the
    pointwise product
    \[
        \chi_{\mathrm{std}\otimes \mathrm{std}}
        =
        \chi_{\mathrm{std}}\chi_{\mathrm{std}}
        =
        (4,0,4,0,0).
    \]
    The multiplicity of \(\mathrm{std}\) in this tensor product is
    \[
        (\chi_{\mathrm{std}},\chi_{\mathrm{std}\otimes \mathrm{std}})
        =
        \frac{1}{8}(2\cdot 4 + (-2)\cdot 4)
        =
        0.
    \]
    For each one-dimensional character \(\chi\), the values on the classes
    represented by \(1\) and \(r^2\) are both \(1\), so
    \[
        (\chi,\chi_{\mathrm{std}\otimes \mathrm{std}})
        =
        \frac{1}{8}(4+4)
        =
        1.
    \]
    Therefore
    \[
        \mathrm{std}\otimes\mathrm{std}
        \cong
        (1,1)\oplus(-1,1)\oplus(1,-1)\oplus(-1,-1).
    \]
}

\corp{
    Let \(G\) be a finite group, and let
    \(V_1,\dots,V_r\) be the irreducible finite-dimensional complex
    representations of \(G\), up to isomorphism. Then the regular
    representation decomposes as
    \[
        \mathbb{C}[G]
        \cong
        \bigoplus_{i=1}^r \dim(V_i)\,V_i.
    \]
    In particular,
    \[
        \sum_{i=1}^r \dim(V_i)^2 = |G|.
    \]
}{
    Let \(\chi_{\mathrm{reg}}\) be the character of the regular representation
    \(\mathbb{C}[G]\). In the basis \(\{e_h\}_{h\in G}\), the action is
    \[
        g\cdot e_h = e_{gh}.
    \]
    If \(g=e\), this is the identity on a vector space of dimension \(|G|\), so
    \[
        \chi_{\mathrm{reg}}(e)=|G|.
    \]
    If \(g\neq e\), then \(h\mapsto gh\) has no fixed points. Hence the matrix
    of \(g\) in the basis \(\{e_h\}_{h\in G}\) is a permutation matrix with no
    diagonal entries equal to \(1\), and therefore
    \[
        \chi_{\mathrm{reg}}(g)=0.
    \]

    By the previous corollary, the multiplicity of \(V_i\) in
    \(\mathbb{C}[G]\) is
    \[
        (\chi_{V_i},\chi_{\mathrm{reg}})
        =
        \frac{1}{|G|}\sum_{g\in G}
        \overline{\chi_{V_i}(g)}\chi_{\mathrm{reg}}(g).
    \]
    Since \(\chi_{\mathrm{reg}}\) vanishes away from \(\mathrm{id}\), this becomes
    \[
        (\chi_{V_i},\chi_{\mathrm{reg}})
        =
        \frac{1}{|G|}
        \overline{\chi_{V_i}(\mathrm{id})}\,|G|
        =
        \dim(V_i).
    \]
    Thus
    \[
        \mathbb{C}[G]
        \cong
        \bigoplus_{i=1}^r \dim(V_i)\,V_i.
    \]
    Taking dimensions on both sides gives
    \[
        |G|
        =
        \dim\mathbb{C}[G]
        =
        \sum_{i=1}^r \dim(V_i)\dim(V_i)
        =
        \sum_{i=1}^r \dim(V_i)^2.
    \]
}{}

\rmkb{
    The previous corollary also gives a useful relation among the irreducible
    characters. Since
    \[
        \chi_{\mathrm{reg}}
        =
        \sum_{i=1}^r \dim(V_i)\chi_{V_i},
    \]
    and since \(\chi_{\mathrm{reg}}(g)=0\) for every \(g\neq \mathrm{id}\), one has
    \[
        0
        =
        \sum_{i=1}^r \dim(V_i)\chi_{V_i}(g)
        \qquad
        \forall g\neq \mathrm{id}.
    \]
    Thus, once all irreducible characters except one are known, the remaining
    one is determined on all non-identity conjugacy classes by this equation.
    Its value at \(\mathrm{id}\) is its dimension, which is constrained by
    \[
        \sum_{i=1}^r \dim(V_i)^2 = |G|.
    \]
}

\thmp{Irreducible Characters Form A Basis}{
    Let \(G\) be a finite group. The characters of the irreducible
    finite-dimensional complex representations of \(G\) form an orthonormal
    basis of the vector space of class functions on \(G\).

    In particular, the number of irreducible representations of \(G\), up to
    isomorphism, is equal to the number of conjugacy classes of \(G\).
}{
    We already know that the irreducible characters are orthonormal. It remains
    to show that they span the space of class functions.

    Let \(\alpha : G \to \mathbb{C}\) be a class function such that
    \[
        (\alpha,\chi_V)=0
    \]
    for every irreducible representation \(V\). We will show that
    \(\alpha=0\).

    For an irreducible representation \((V,\rho)\), define
    \[
        \varphi_{\alpha,V}
        :=
        \sum_{g\in G}\alpha(g)\rho(g)
        \in \operatorname{End}(V).
    \]
    Since \(\alpha\) is a class function, \(\varphi_{\alpha,V}\) commutes with the
    \(G\)-action. Indeed, for \(h\in G\),
    \[
        \varphi_{\alpha,V}\rho(h)
        =
        \sum_{g\in G}\alpha(g)\rho(gh).
    \]
    Writing \(g'=h^{-1}gh\), this becomes
    \[
        \varphi_{\alpha,V}\rho(h)
        =
        \sum_{g'\in G}\alpha(hg'h^{-1})\rho(hg')
        =
        \rho(h)\sum_{g'\in G}\alpha(g')\rho(g')
        =
        \rho(h)\varphi_{\alpha,V}.
    \]
    Thus \(\varphi_{\alpha,V}\) is a map of \(G\)-representations
    \(V\to V\). By Schur's lemma,
    \[
        \varphi_{\alpha,V} = \lambda_V 1_V
    \]
    for some \(\lambda_V\in\mathbb{C}\).

    Taking traces gives
    \[
        \lambda_V
        =
        \frac{1}{\dim(V)}\operatorname{tr}(\varphi_{\alpha,V})
        =
        \frac{1}{\dim(V)}\sum_{g\in G}\alpha(g)\chi_V(g)
        =
        \frac{|G|}{\dim(V)}
        \overline{(\alpha,\chi_{V^*})}
        =
        0.
    \]
    Here \(V^*\) is again irreducible, so the last equality follows from the
    assumption on \(\alpha\). Hence \(\varphi_{\alpha,V}=0\) for every
    irreducible representation \(V\).

    Therefore, for any finite-dimensional representation \((W,\rho_W)\), one
    has
    \[
        \sum_{g\in G}\alpha(g)\rho_W(g)=0,
    \]
    because \(W\) is a direct sum of irreducibles.

    Apply this to the regular representation \(\mathbb{C}[G]\). The operators
    \(\rho_{\mathrm{reg}}(g)\) are linearly independent: indeed,
    \[
        \rho_{\mathrm{reg}}(g)(e_{\mathrm{id}})=e_g.
    \]
    Hence
    \[
        \alpha(g)=0
        \qquad
        \forall g\in G.
    \]
    Thus \(\alpha=0\), and the irreducible characters span the class functions.

    Since the irreducible characters form a basis, their number is the
    dimension of the space of class functions. A class function is determined
    by its values on conjugacy classes, so this dimension is the number of
    conjugacy classes of \(G\).
}

\section{Representations Of \texorpdfstring{\(S_n\)}{Sn}}

The irreducible representations of \(S_n\) are indexed by the conjugacy
classes of \(S_n\), and hence by the partitions of \(n\).

\defn{Partitions And Young Diagrams}{
    A \textbf{partition} of \(n\) is a tuple
    \[
        \lambda = (\lambda_1,\dots,\lambda_k)
    \]
    such that
    \[
        n = \sum_{i=1}^k \lambda_i,
        \qquad
        \lambda_1 \geq \lambda_2 \geq \cdots \geq \lambda_k > 0.
    \]
    We write \(\lambda \vdash n\).

    A partition can be represented by a \textbf{Young diagram}: draw
    \(\lambda_i\) boxes in the \(i\)-th row, aligned on the left.
}

\ex{
    The partitions of \(3\) are
    \[
        (3),
        \qquad
        (2,1),
        \qquad
        (1,1,1).
    \]
    Their Young diagrams are
    \[
        \ydiagram{3}
        \qquad
        \ydiagram{2,1}
        \qquad
        \ydiagram{1,1,1}.
    \]
    Thus \(S_3\) has three irreducible representations, corresponding to these
    three diagrams.
}

\defn{Associative Algebras}{
    An \textbf{associative algebra} \((A,\mu,1)\) is a vector space \(A\)
    equipped with a linear map
    \[
        \mu : A \otimes A \to A,
        \qquad
        a\otimes b \mapsto ab,
    \]
    and an element \(1\in A\) such that
    \[
        (ab)c = a(bc)
        \qquad
        \forall a,b,c\in A,
    \]
    and
    \[
        1a=a=a1
        \qquad
        \forall a\in A.
    \]
}

\ex{
    Let \(G\) be a group. The vector space \(\mathbb{C}[G]\), with basis
    \[
        \{e_g\}_{g\in G},
    \]
    is an associative algebra with multiplication
    \[
        e_g e_h = e_{gh},
    \]
    and unit
    \[
        1=e_{\mathrm{id}}.
    \]
    It is called the \textbf{group algebra} of \(G\).
}

\rmkb{
    A representation
    \[
        \rho : G \to \operatorname{Aut}(V)
    \]
    induces a linear map
    \[
        \widehat{\rho} : \mathbb{C}[G] \to \operatorname{End}(V),
        \qquad
        e_g \mapsto \rho(g).
    \]
    In particular,
    \[
        \widehat{\rho}\left(\frac{1}{|G|}\sum_{g\in G} e_g\right)
    \]
    is the projection onto the trivial component of \(V\).
}

\defn{Tableaux}{
    A \textbf{tableau} \(T\) on a given Young diagram with \(n\) boxes is a
    numbering of its boxes by \(1,\dots,n\).

    Associated to \(T\) are two subgroups of \(S_n\):
    \[
        P_T
        :=
        \{\sigma\in S_n \mid \sigma \text{ preserves the rows of } T\},
    \]
    and
    \[
        Q_T
        :=
        \{\sigma\in S_n \mid \sigma \text{ preserves the columns of } T\}.
    \]
}

\ex{
    For example, for the tableau
    \[
        T=
        \begin{ytableau}
            1 & 2 & 3 \\
            4 & 5 \\
            6 & 7
        \end{ytableau}
    \]
    one has
    \[
        P_T \cong S_3\times S_2\times S_2,
        \qquad
        Q_T \cong S_3\times S_3\times S_1.
    \]

    We will suppress the tableau from the notation and write \(\lambda\) for
    both a partition and a choice of tableau of shape \(\lambda\).
}

\defn{Young Symmetrizers}{
    Let \(\lambda\vdash n\), and choose a tableau of shape \(\lambda\). Define
    \[
        a_\lambda
        :=
        \sum_{g\in P_\lambda} e_g
        \in \mathbb{C}[S_n],
    \]
    and
    \[
        b_\lambda
        :=
        \sum_{g\in Q_\lambda}\operatorname{sign}(g)e_g
        \in \mathbb{C}[S_n].
    \]
    The \textbf{Young symmetrizer} associated to \(\lambda\) is
    \[
        c_\lambda
        :=
        a_\lambda b_\lambda
        \in \mathbb{C}[S_n].
    \]
}

\ex{
    For \(\lambda=(n)\), the Young diagram has one row. Hence
    \[
        P_\lambda=S_n,
        \qquad
        Q_\lambda=\{\mathrm{id}\},
    \]
    and so
    \[
        c_\lambda=a_\lambda.
    \]
    In this case
    \[
        \operatorname{im}((\cdot)c_\lambda)\cong \mathrm{Triv}.
    \]

    For \(\lambda=(1,\dots,1)\), the Young diagram has one column. Hence
    \[
        P_\lambda=\{\mathrm{id}\},
        \qquad
        Q_\lambda=S_n,
    \]
    and so
    \[
        a_\lambda=1,
        \qquad
        c_\lambda=b_\lambda.
    \]
    In this case
    \[
        \operatorname{im}((\cdot)c_\lambda)\cong \mathrm{Alt}.
    \]
}

\defn{Representations From Young Symmetrizers}{
    Let \(\lambda\vdash n\). Define
    \[
        V_\lambda
        :=
        \mathbb{C}[S_n]c_\lambda
        =
        \{a c_\lambda \mid a\in \mathbb{C}[S_n]\}
        \leq \mathbb{C}[S_n].
    \]
    This is a representation of \(S_n\) by left multiplication:
    \[
        \rho(g) : V_\lambda \to V_\lambda,
        \qquad
        a c_\lambda \mapsto e_g a c_\lambda.
    \]
}

\ex{
    For \(S_3\), the three partitions give
    \[
        V_{\ydiagram{3}}
        =
        \mathbb{C}\bigl[
            e_{\mathrm{id}}+e_{(12)}+e_{(23)}+e_{(13)}
            +e_{(123)}+e_{(132)}
        \bigr],
    \]
    and
    \[
        V_{\ydiagram{1,1,1}}
        =
        \mathbb{C}\bigl[
            e_{\mathrm{id}}-e_{(12)}-e_{(23)}-e_{(13)}
            +e_{(123)}+e_{(132)}
        \bigr].
    \]
    For the tableau
    \[
        \begin{ytableau}
            1 & 2 \\
            3
        \end{ytableau}
    \]
    one has
    \[
        c_\lambda
        =
        (e_{\mathrm{id}}+e_{(12)})(e_{\mathrm{id}}-e_{(13)})
        =
        e_{\mathrm{id}}+e_{(12)}-e_{(13)}-e_{(123)}.
    \]
    Hence
    \[
        V_{\ydiagram{2,1}}
        =
        \mathbb{C}[S_3]c_\lambda.
    \]
}

\prop{
    Let \(T,T'\) be Young tableaux of the same shape \(\lambda\). Then there
    exists a unique \(g\in S_n\) such that
    \[
        T'=gT.
    \]
    If \(P,Q\) are the row and column subgroups associated to \(T\), and
    \(P',Q'\) are the row and column subgroups associated to \(T'\), then
    \[
        P'=gPg^{-1},
        \qquad
        Q'=gQg^{-1}.
    \]
    Hence
    \[
        c_{\lambda,T'}
        =
        e_g c_{\lambda,T} e_{g^{-1}}.
    \]
    Therefore
    \[
        V_{\lambda,T}\to V_{\lambda,T'},
        \qquad
        a c_{\lambda,T}\mapsto a c_{\lambda,T}e_{g^{-1}}
        =
        a e_{g^{-1}} c_{\lambda,T'},
    \]
    is an isomorphism of \(S_n\)-representations.

    Moreover,
    \[
        P\cap Q=\{\mathrm{id}\}.
    \]
    Thus, if
    \[
        g=pq=p'q',
        \qquad
        p,p'\in P,
        \quad
        q,q'\in Q,
    \]
    then
    \[
        p'^{-1}p=qq'^{-1}\in P\cap Q.
    \]
    Hence \(p=p'\) and \(q=q'\), so the decomposition \(g=pq\) is unique.
    Consequently,
    \[
        c_{\lambda,T}
        =
        \sum_{g=pq}\operatorname{sign}(q)e_g.
    \]
}

\propp{
    For each Young tableau, the representation \(V_\lambda\) is irreducible.
}{
    Let \(v\in V_\lambda\). We want to understand \(c_\lambda v\). Write
    \[
        v=a c_\lambda.
    \]
    Then \(c_\lambda v=c_\lambda a c_\lambda\). Hence \(c_\lambda v\) solves
    the equation
    \[
        e_p x e_q\operatorname{sign}(q)=x
        \qquad
        \forall p\in P,\ q\in Q.
        \tag{1}
    \]

    Write
    \[
        x=\sum_{g\in S_n}n_g e_g.
    \]
    Equation (1) implies
    \[
        n_{pgq}=n_g\operatorname{sign}(q).
    \]
    If \(g\neq pq\), then let \(T'=gT\). Assume that there are integers
    \(i<j\leq n\) such that \(i,j\) are in the same row of \(T\) and in the
    same column of \(T'\). Then \(t=(ij)\in P\), and also
    \[
        t\in Q'=gQg^{-1}.
    \]
    Thus there is some \(q\in Q\) with \(t=gqg^{-1}\), and hence
    \[
        g=tgq^{-1}.
    \]
    Comparing coefficients gives
    \[
        n_g=-n_g,
        \qquad
        \text{so } n_g=0.
    \]

    Now let \(T,T'=gT\) be two Young tableaux such that no pair \(i,j\) is in
    the same row of \(T\) and in the same column of \(T'\). Then there are
    \(p\in P\) and \(q'\in Q'\) such that the first rows of \(pT\) and
    \(q'T'\) agree. Continuing row by row, we get
    \[
        pT=q'T'=q'gT.
    \]
    Hence
    \[
        p=q'g.
    \]
    Since \(q'\in Q'=gQg^{-1}\), we can write \(q'=gqg^{-1}\) for some
    \(q\in Q\). Therefore
    \[
        p=gq,
        \qquad
        g=pq^{-1}.
    \]
    It follows that the only possible non-zero coefficients are those with
    \(g=pq\). Thus every solution of (1) is of the form
    \[
        x=n_1\sum_{g=pq}e_g\operatorname{sign}(q)
        =
        n_1 c_\lambda.
    \]
    Hence
    \[
        c_\lambda V_\lambda\subseteq \mathbb{C}[c_\lambda].
    \]

    Let \(W\subseteq V_\lambda\) be a subrepresentation. Then either
    \[
        c_\lambda W=0
        \qquad\text{or}\qquad
        c_\lambda W=\mathbb{C}[c_\lambda].
    \]
    In the second case, \(c_\lambda\in W\), and hence
    \[
        a c_\lambda\in W
        \qquad
        \forall a\in \mathbb{C}[S_n].
    \]
    Thus \(V_\lambda\subseteq W\), so \(W=V_\lambda\).

    In the first case, choose a projection operator
    \[
        P_W:\mathbb{C}[S_n]\to \mathbb{C}[S_n]
    \]
    which projects onto \(W\). It satisfies
    \[
        e_gP_W(a)=P_W(e_ga).
    \]
    For \(a=1\), this gives
    \[
        e_gP_W(1)=P_W(e_g).
    \]
    Hence \(P_W\) is right multiplication by \(P_W(1)\). Since
    \(P_W(1)\in W\), the assumption \(c_\lambda W=0\) gives
    \[
        P_W(c_\lambda)=c_\lambda P_W(1)=0.
    \]
    Therefore, for every \(a\in\mathbb{C}[S_n]\),
    \[
        P_W(a c_\lambda)=aP_W(c_\lambda)=0.
    \]
    Since \(W\subseteq V_\lambda\) and \(P_W\) is the identity on \(W\), we get
    \(W=0\).

    Thus the only subrepresentations are \(0\) and \(V_\lambda\).
}

\rmkb{
    The proof shows that
    \[
        c_\lambda V_\lambda\neq 0.
    \]
}

\thmp{Classification Of Irreducible Representations Of \texorpdfstring{\(S_n\)}{Sn}}{
    For \(\lambda\neq\mu\), the representations \(V_\lambda\) and \(V_\mu\)
    are not isomorphic.

    Hence the \(V_\lambda\), where \(\lambda\vdash n\), form a complete set of
    irreducible representations of \(S_n\).
}{
    Assume \(\lambda>\mu\), meaning that the first non-zero difference
    \(\lambda_i-\mu_i\) is positive. We show that
    \[
        c_\lambda V_\mu=\{0\},
        \qquad
        c_\lambda V_\lambda\neq 0.
    \]
    The second statement was shown in the previous proof. For the first one,
    it is enough to show that
    \[
        c_\lambda a c_\mu=0
        \qquad
        \forall a\in\mathbb{C}[S_n].
    \]
    Since
    \[
        c_\lambda a c_\mu
        =
        a_\lambda b_\lambda a a_\mu b_\mu,
    \]
    it is enough to prove
    \[
        a_\lambda x b_\mu=0
        \qquad
        \forall x\in\mathbb{C}[S_n].
    \]

    It suffices to take \(x=e_g\). Then
    \[
        a_\lambda e_g b_\mu
        =
        a_\lambda b_{gT_\mu}.
    \]
    Since \(\lambda>\mu\), there exist \(i,j\) which lie in the same row of
    \(T_\lambda\) and in the same column of \(gT_\mu\). Let \(t=(ij)\). Then
    \(t\in P_\lambda\), so
    \[
        a_\lambda e_t=a_\lambda,
    \]
    while \(t\) lies in the column group of \(gT_\mu\), so
    \[
        e_t b_{gT_\mu}=-b_{gT_\mu}.
    \]
    Therefore
    \[
        a_\lambda b_{gT_\mu}
        =
        a_\lambda e_t e_t b_{gT_\mu}
        =
        -a_\lambda b_{gT_\mu}.
    \]
    Hence
    \[
        a_\lambda b_{gT_\mu}=0.
    \]
    Thus \(c_\lambda V_\mu=\{0\}\).

    Since \(c_\lambda V_\lambda\neq 0\), the representations
    \(V_\lambda\) and \(V_\mu\) cannot be isomorphic. We have already shown
    that every \(V_\lambda\) is irreducible, and there are as many partitions
    \(\lambda\vdash n\) as conjugacy classes in \(S_n\). Therefore these
    representations classify all irreducible representations of \(S_n\).
}

\subsection{Frobenius Formula}

For
\[
    i=(i_1,\dots,i_n),
    \qquad
    \sum_{\alpha=1}^{n} i_\alpha \alpha=n,
\]
we denote by \(C_i\) the conjugacy class of \(S_n\) consisting of
permutations with \(i_\alpha\) many \(\alpha\)-cycles.

Let \(\lambda\vdash n\) be a partition with \(k\) rows. We work with
polynomials in \(k\) variables
\[
    x_1,\dots,x_k.
\]
For \(1\leq j\leq n\), define
\[
    p_j(x)
    :=
    x_1^j+\cdots+x_k^j,
\]
and define
\[
    \Delta(x)
    :=
    \prod_{a<b}(x_a-x_b).
\]
For a polynomial \(Q(x)\) in \(x_1,\dots,x_k\), we denote by
\[
    [Q(x)]_{\ell_1,\dots,\ell_k}
\]
the coefficient of
\[
    x_1^{\ell_1}\cdots x_k^{\ell_k}.
\]
Here
\[
    \ell_1=\lambda_1+k-1,
    \qquad
    \ell_2=\lambda_2+k-2,
    \qquad
    \dots,
    \qquad
    \ell_k=\lambda_k.
\]

\thmp{Frobenius Formula}{
    With the notation above,
    \[
        \chi_\lambda(C_i)
        =
        \left[
            \Delta(x)
            \prod_{j=1}^{n} p_j(x)^{i_j}
        \right]_{\ell_1,\dots,\ell_k}.
    \]
}{
    Look it up.
}

\ex{
    Let \(n=5\), and let
    \[
        \lambda=(3,2)
        =
        \ydiagram{3,2}.
    \]
    Consider
    \[
        i=(0,1,1,0,0),
    \]
    so that
    \[
        (12)(345)\in C_i.
    \]
    Here \(k=2\), hence
    \[
        \ell_1=\lambda_1+1=4,
        \qquad
        \ell_2=\lambda_2=2.
    \]
    Therefore the Frobenius formula gives
    \[
        \chi_\lambda(C_i)
        =
        \left[
            (x_1-x_2)(x_1^2+x_2^2)(x_1^3+x_2^3)
        \right]_{4,2}
        =
        1.
    \]
}

\rmkb{
    The representation corresponding to
    \[
        \ydiagram{4,1}
    \]
    is the standard representation: \(S_n\) acts on \(\mathbb{C}^n\) by
    permuting coordinates. This representation is not irreducible, since
    \[
        e_1+\cdots+e_n
        \in
        (\mathbb{C}^n)^{S_n}.
    \]
    The subspace
    \[
        V
        :=
        \left\{
            (x_1,\dots,x_n)\in\mathbb{C}^n
            \,\middle|\,
            \sum_{i=1}^{n}x_i=0
        \right\}
    \]
    is the standard \(S_n\)-representation.

    More generally, the representation corresponding to
    \[
        (n-k,1,\dots,1)
    \]
    is
    \[
        \bigwedge^k V.
    \]
    If \(\lambda'\) denotes the conjugate partition, then
    \[
        V_\lambda
        =
        V_{\lambda'}\otimes\mathrm{Alt}.
    \]
}
